\documentclass[11pt]{article}
\usepackage[margin=1in]{geometry}
\usepackage{xcolor}
\usepackage{enumitem}
\usepackage{amsmath}
\usepackage{amssymb}
\usepackage{hyperref}
\hypersetup{colorlinks=true,linkcolor=black,urlcolor=blue}
\definecolor{revblue}{RGB}{0,0,180}
\newcommand{\loc}[1]{\textbf{[\textcolor{revblue}{Tracked PDF: #1}]}}
\setlist{nosep,leftmargin=1.4em}
\setlength{\parskip}{0.55em}
\setlength{\parindent}{0pt}

\begin{document}

\begin{center}
{\Large\textbf{Cover Letter to Reviewer 3}}\\[0.5em]
\textbf{Round 2 Revision}
\end{center}

\textbf{Manuscript:} ``Content-Conditioned Invariant Feature-Statistic
Augmentation: Unifying Robustness and Weight Averaging for Domain
Generalization'' (CIFA)\\
\textbf{Journal:} \textit{Electronics} (MDPI)\qquad
\textbf{Manuscript ID:} electronics-4654208\\
\textbf{Date:} 3 November 2026

\hrule
\vspace{0.4em}

Dear Reviewer 3,

We sincerely thank you for the detailed and constructive assessment and for
recognising the substantial improvement in the mathematical formulation,
experimental validation, statistical analysis, and interpretation. We have
addressed the three limited revisions. Each point below is preceded by a
location tag giving the exact page and line range in the tracked PDF, which
uses continuous line numbering.

\subsection*{Point 1 (R3-1) --- Mathematical justification for the conditional posterior bound}
\loc{Page 8, Line 304 through Page 9, Line 348; proofs on Page 10, Line 354
through Page 11, Line 400}

Your comment is correct, and we have restructured Proposition~1 rather than
defend the earlier statement. The data-processing inequality controls the
distributions (laws) of the posterior random variables, not the posterior
evaluated at each fixed observation; at rare observations a small likelihood
perturbation can be amplified. The proposition is now stated at two clearly
separated levels:

\begin{itemize}
\item \textbf{(L2-i) Exact case.} When $\tau=0$ and the partner source is
exogenous, the content posterior is preserved exactly;
\item \textbf{(L2-ii)(a) Distribution-level bounds (unconditional).} The
coupling gives the only unconditional statements:
\[
\mathrm{TV}(P_{C,Z},P_{C,T(Z)})\le\tau,\qquad
\mathrm{TV}(p(Z),p(T(Z)))\le\tau;
\]
\item \textbf{(L2-ii)(b) Pointwise bound (conditional).} Retained only under
an explicit inter-model likelihood-ratio condition
$|\log p(\mathbf{b},r,\mathbf{d}\mid c)-\log
p(\mathbf{b},r,\mathbf{d}_{aug}\mid c)|\le\alpha$, under which
$\mathrm{TV}(p(C\mid Z),p(C\mid T(Z)))\le\tanh\alpha$.
\end{itemize}

The KL corollary is now conditional on the same assumption,
$D_{\mathrm{KL}}\le\tanh(\alpha)/q_{\min}$, and the earlier
$\tau\log(M/p_{\min})$ and $\tau/q_{\min}$ pointwise bounds are removed. The
label-posterior drift $\delta$ is treated as a directly measured quantity
rather than added to $\tau$. The full derivations are given in the revised
proof (Pages~10--11). Theorems~1 and~2 are unaffected because they use the
distribution-level total-variation quantity.

\subsection*{Point 2 (R3-2) --- Consistent distinction of equivalence, non-significance, and superiority}
\loc{Page 1, Lines 17--19; Pages 21--23, Lines 795--828; Page 36,
Lines 1168--1200}

\begin{sloppypar}
We have applied one consistent terminology throughout the twelve affected
passages. Non-significant contrasts are reported as statistically tied or on
par; the term equivalence is used only for the TOST procedure and always with
the equivalence margin stated; and no superiority is claimed. We also corrected
the TOST implementation to the standard two-one-sided procedure using 90\%
confidence intervals and recomputed every contrast. On PACS ($n=10$), at the
one-point margin equivalence is confirmed only against RSC; at the two-point
margin it additionally holds against SWAD, CORAL, and Fish; the CIFA--EFDM
contrast is neither superior nor certified equivalent. These changes appear at
the abstract (Page~1), introduction and contributions (Pages~2--3),
experiments and captions (Pages~21--25), appendix (Pages~33--35), and
conclusions (Page~36). The conditional nature of the unification is kept
explicit throughout.
\end{sloppypar}

\subsection*{Point 3 (R3-3) --- Clean manuscript and consistency check}
\loc{Whole manuscript; hyphenation at Page 36, Line 1168}

A clean, fully accepted version (\texttt{CIFA\_Electronics\_clean.pdf},
39 pages) is provided alongside the tracked version. We have verified that
every cross-reference resolves with no dangling or duplicate labels; all 27
bibliography entries are cited; equations, tables, and figures use automatic
numbering with no hard-coded numbers; and the terminology and notation are
consistent. The clean manuscript renders entirely in black with no revision
marks. As part of this pass we also standardised the hyphenation (e.g.,
``content--style''). Both the clean and tracked manuscripts compile with zero
errors and zero overfull boxes, and their accepted content is identical.

\vspace{0.6em}
We thank you once again for the constructive and supportive review, and hope
these revisions fully address the remaining points.

\vspace{0.8em}
Sincerely,\\
Qingsong Cai (on behalf of all authors)\\
\textit{Beijing Technology and Business University}\\
\textit{caiqs@btbu.edu.cn}

\end{document}
